\chapter{Chapter09 — Small Set Expansion}
%%% generated by the blueprint pipeline from FABL.Chapter09.SmallSetExpansion
\section{Overview}

Declarations of the Lean module \texttt{FABL.Chapter09.SmallSetExpansion}.

\section{Declarations}

\begin{definition}[Two-four noise parameter]
\lean{FABL.twoFourNoiseRoot}
\label{FABL.twoFourNoiseRoot}
\leanok
Define the two-four noise parameter to be $\sqrt{1/3}$.
\end{definition}

\begin{theorem}[Nonnegativity of the two-four noise root]
\lean{FABL.twoFourNoiseRoot_nonneg}
\label{FABL.twoFourNoiseRoot_nonneg}
\leanok
\uses{FABL.twoFourNoiseRoot}
\difficulty{1}
The number $\sqrt{1/3}$ is nonnegative.
\end{theorem}

\begin{theorem}[Square of the two-four noise root]
\lean{FABL.twoFourNoiseRoot_sq}
\label{FABL.twoFourNoiseRoot_sq}
\leanok
\uses{FABL.twoFourNoiseRoot}
\difficulty{2}
Let $r = \sqrt{1/3}$. Then $r^2 = 1/3$.
\end{theorem}

\begin{theorem}[Noise parameter lies in the unit interval]
\lean{FABL.twoFourNoiseRoot_mem_Icc}
\label{FABL.twoFourNoiseRoot_mem_Icc}
\leanok
\uses{FABL.twoFourNoiseRoot, FABL.twoFourNoiseRoot_nonneg}
\difficulty{4}
The number $\sqrt{1/3}$ belongs to the interval $[-1,1]$.
\end{theorem}

\begin{theorem}[Uniqueness from Fourier coefficients]
\lean{FABL.smallSet_fourier_ext}
\label{FABL.smallSet_fourier_ext}
\leanok
\uses{FABL.fourierCoeff, FABL.fourier_expansion}
\difficulty{4}
Let $n$ be a natural number, and let $f$ and $g$ be real-valued functions on the
$n$-dimensional Boolean cube. If their Fourier coefficients agree for every subset of
the coordinate set, then $f=g$.
\end{theorem}

\begin{definition}[Odd part in the first coordinate]
\lean{FABL.hypercontractiveOddPart}
\label{FABL.hypercontractiveOddPart}
\leanok
\uses{FABL.firstCoordinateSlice}
For a function $f : \{−1,1\}^{n+1} \to \mathbb{R}$ and a point $x \in \{−1,1\}^n$,
define its odd part in the first coordinate to be half the difference between the values
of $f$ at the two points whose first coordinates are $1$ and $-1$, respectively, and
whose remaining coordinates are $x$.
\end{definition}

\begin{definition}[Even part in the first coordinate]
\lean{FABL.hypercontractiveEvenPart}
\label{FABL.hypercontractiveEvenPart}
\leanok
\uses{FABL.firstCoordinateSlice}
For a function $f : \{−1,1\}^{n+1} \to \mathbb{R}$ and a point $x \in \{−1,1\}^n$,
define its even part in the first coordinate to be half the sum of the values of $f$ at
the two points whose first coordinates are $1$ and $-1$, respectively, and whose
remaining coordinates are $x$.
\end{definition}

\begin{theorem}[Fourier coefficients of the odd slice]
\lean{FABL.fourierCoeff_hypercontractiveOddPart}
\label{FABL.fourierCoeff_hypercontractiveOddPart}
\leanok
\uses{FABL.firstCoordinateSlice, FABL.fourierCoeff, FABL.fourierCoeff_insert_zero_tailFrequency, FABL.hypercontractiveOddPart, FABL.monomial, FABL.tailFrequency}
\difficulty{4}
Let $n$ be a natural number, let $f$ be a real-valued function on the sign cube of
dimension $n+1$, and let $S$ be a subset of the first $n$ coordinate indices. Take half
the difference between the slices of $f$ obtained by fixing its first coordinate to $1$
and $-1$. The Fourier coefficient of this function at $S$ equals the Fourier coefficient
of $f$ at the frequency consisting of the first coordinate together with the indices of
$S$ shifted up by one.
\end{theorem}

\begin{theorem}[Fourier coefficients of the even slice]
\lean{FABL.fourierCoeff_hypercontractiveEvenPart}
\label{FABL.fourierCoeff_hypercontractiveEvenPart}
\leanok
\uses{FABL.firstCoordinateSlice, FABL.fourierCoeff, FABL.fourierCoeff_tailFrequency, FABL.hypercontractiveEvenPart, FABL.monomial, FABL.tailFrequency}
\difficulty{4}
Let $n$ be a natural number, let $f$ be a real-valued function on the sign cube of
dimension $n+1$, and let $S$ be a subset of the first $n$ coordinate indices. Average
the two slices of $f$ obtained by fixing its first coordinate to $1$ and $-1$. The
Fourier coefficient of this averaged function at $S$ equals the Fourier coefficient of
$f$ at the frequency obtained by shifting every index in $S$ up by one.
\end{theorem}

\begin{theorem}[Fourier coefficient of a coordinate slice]
\lean{FABL.fourierCoeff_slice_tail_add_insert}
\label{FABL.fourierCoeff_slice_tail_add_insert}
\leanok
\uses{FABL.Sign, FABL.firstCoordinateSlice, FABL.fourierCoeff, FABL.fourierCoeff_insert_zero_tailFrequency, FABL.fourierCoeff_tailFrequency, FABL.signValue, FABL.signValue_neg_one, FABL.signValue_one, FABL.tailFrequency}
\difficulty{4}
Let $n$ be a natural number, let $f:\{-1,1\}^{n+1}\to\mathbb R$, let
$b\in\{-1,1\}$, and let $S\subseteq\{1,\ldots,n\}$. Write
$S^+\subseteq\{1,\ldots,n+1\}$ for the set obtained by shifting every coordinate in
$S$ up by one. The Fourier coefficient of the slice of $f$ with first coordinate fixed
to $b$, at $S$, equals the Fourier coefficient of $f$ at $S^+$ plus $b$ times its
coefficient at $S^+\cup\{1\}$.
\end{theorem}

\begin{theorem}[First-coordinate slice of the noise operator]
\lean{FABL.firstCoordinateSlice_noiseOperator}
\label{FABL.firstCoordinateSlice_noiseOperator}
\leanok
\uses{FABL.Sign, FABL.card_tailFrequency, FABL.firstCoordinateSlice, FABL.fourierCoeff, FABL.fourierCoeff_hypercontractiveEvenPart, FABL.fourierCoeff_hypercontractiveOddPart, FABL.fourierCoeff_noiseOperator, FABL.fourierCoeff_slice_tail_add_insert, FABL.hypercontractiveEvenPart, FABL.hypercontractiveOddPart, FABL.monomial, FABL.noiseOperator, FABL.signValue, FABL.smallSet_fourier_ext, FABL.tailFrequency, FABL.zero_notMem_tailFrequency}
\difficulty{4}
Let $n\in\mathbb N$, $\rho\in\mathbb R$, let $f$ be a real-valued function on the
$(n+1)$-dimensional sign cube, and let $b\in\{\pm1\}$. The slice at first coordinate $b$
of the noise operator applied to $f$ equals, as a function of $x\in\{\pm1\}^n$, $b\rho$
times the noise operator applied to the odd part of $f$, plus the noise operator applied
to its even part.
\end{theorem}

\begin{theorem}[Nonnegative second moment]
\lean{FABL.smallSet_expect_sq_nonneg}
\label{FABL.smallSet_expect_sq_nonneg}
\leanok
\difficulty{1}
For every natural number $n$ and every real-valued function $f$ on the $n$-dimensional
Boolean cube, the uniform expectation of $f(x)^2$ is nonnegative.
\end{theorem}

\begin{theorem}[Nonnegative fourth moment]
\lean{FABL.smallSet_expect_fourth_nonneg}
\label{FABL.smallSet_expect_fourth_nonneg}
\leanok
\difficulty{1}
For every natural number $n$ and every real-valued function $f$ on the $n$-dimensional
Boolean cube, the uniform expectation of $f(x)^4$ is nonnegative.
\end{theorem}

\begin{theorem}[Nonnegative product of squared functions]
\lean{FABL.smallSet_expect_sq_mul_sq_nonneg}
\label{FABL.smallSet_expect_sq_mul_sq_nonneg}
\leanok
\difficulty{1}
For every natural number $n$ and all real-valued functions $d$ and $e$ on the
$n$-dimensional Boolean cube, the uniform expectation of $d(x)^2e(x)^2$ is nonnegative.
\end{theorem}

\begin{theorem}[Fourth moment bound under slice hypotheses]
\lean{FABL.fourthMoment_le_sq_of_hypercontractiveSlice}
\proofsource{Analysis-of-Boolean-Functions-by-Ryan-ODonnell}{pdf-ffbe9c7b3306-p253-b003}
\label{FABL.fourthMoment_le_sq_of_hypercontractiveSlice}
\leanok
\uses{FABL.expect_fourth_eq_expect_odd_even, FABL.expect_sq_mul_sq_sq_le, FABL.signValue, FABL.smallSet_expect_fourth_nonneg, FABL.smallSet_expect_sq_mul_sq_nonneg}
\difficulty{5}
Let $n$ be a nonnegative integer and let $d,e:\{-1,1\}^n\to\mathbb R$ and
$f:\{-1,1\}^{n+1}\to\mathbb R$. Suppose that $f(b,x)=bd(x)+e(x)$ for every
$b\in\{-1,1\}$ and $x\in\{-1,1\}^n$. Let $A,B\in\mathbb R$ satisfy $0\le A$ and $0\le
B$, and suppose that $\mathbb E[d^4]\le A^2/9$ and $\mathbb E[e^4]\le B^2$, with
expectations uniform on the respective cubes. Then $\mathbb E[f^4]\le(A+B)^2$, with
expectation uniform on the $(n+1)$-dimensional cube.
\end{theorem}

\begin{theorem}[Fourth moment hypercontractive inequality]
\lean{FABL.twoFourHypercontractive_fourthMoment}
\proofsource{Analysis-of-Boolean-Functions-by-Ryan-ODonnell}{
  pdf-ffbe9c7b3306-p253-b002,
  pdf-ffbe9c7b3306-p253-b003
}
\label{FABL.twoFourHypercontractive_fourthMoment}
\leanok
\uses{FABL.Sign, FABL.expect_sq_eq_expect_odd_even, FABL.firstCoordinateSlice, FABL.firstCoordinateSlice_noiseOperator, FABL.fourierCoeff_noiseOperator, FABL.fourthMoment_le_sq_of_hypercontractiveSlice, FABL.hypercontractiveEvenPart, FABL.hypercontractiveOddPart, FABL.noiseOperator, FABL.signValue, FABL.smallSet_expect_sq_nonneg, FABL.smallSet_fourier_ext, FABL.twoFourNoiseRoot, FABL.twoFourNoiseRoot_sq}
\difficulty{6}
Let $n$ be a natural number and let $f:\{-1,1\}^n\to\mathbb R$. Applying the noise
operator with parameter $1/\sqrt{3}$, the uniform fourth moment of the resulting
function is at most the square of the uniform second moment of $f$:
\[\mathbb E_x\big[(T_{1/\sqrt{3}}f(x))^4\big]\leq \big(\mathbb E_x[f(x)^2]\big)^2,\]
where $x$ is uniform on $\{-1,1\}^n$.
\end{theorem}

\begin{theorem}[Fourth power of the fourth norm]
\lean{FABL.uniformLpNorm_four_pow_eq_expect_fourth}
\label{FABL.uniformLpNorm_four_pow_eq_expect_fourth}
\leanok
\uses{FABL.uniformLpNorm}
\difficulty{4}
Let $\Omega$ be a finite set and let $f:\Omega\to\mathbb R$. The fourth power of the
uniform $L^4$ norm of $f$ equals the uniform mean of the function $x\mapsto f(x)^4$.
\end{theorem}

\begin{theorem}[Nonnegativity of the uniform norm]
\lean{FABL.uniformLpNorm_nonneg}
\label{FABL.uniformLpNorm_nonneg}
\leanok
\uses{FABL.uniformLpNorm}
\difficulty{1}
Let $\Omega$ be a finite set, let $p\in\mathbb R$, and let $f\colon\Omega\to\mathbb R$.
The uniform $L^p$ norm of $f$ is nonnegative.
\end{theorem}

\begin{theorem}[Two-to-four hypercontractivity]
\lean{FABL.twoFourHypercontractivity}
\proofsource{Analysis-of-Boolean-Functions-by-Ryan-ODonnell}{
  pdf-ffbe9c7b3306-p253-b002,
  pdf-ffbe9c7b3306-p253-b003
}
\label{FABL.twoFourHypercontractivity}
\leanok
\uses{FABL.noiseOperator, FABL.twoFourHypercontractive_fourthMoment, FABL.twoFourNoiseRoot, FABL.uniformLpNorm, FABL.uniformLpNorm_four_pow_eq_expect_fourth, FABL.uniformLpNorm_nonneg, FABL.uniformLpNorm_two_sq_eq_expect_sq}
\difficulty{3}
Let $n$ be a natural number and let $f:\{-1,1\}^n\to\mathbb R$. The uniform $L^4$ norm
of the noise operator applied to $f$ with parameter $1/\sqrt{3}$ is at most the uniform
$L^2$ norm of $f$:
\[\|T_{1/\sqrt{3}}f\|_4\leq \|f\|_2.\]
\end{theorem}

\begin{theorem}[Uniform Hölder inequality]
\lean{FABL.uniformInner_le_uniformLpNorm_fourThird_mul_four}
\label{FABL.uniformInner_le_uniformLpNorm_fourThird_mul_four}
\leanok
\uses{FABL.uniformInner, FABL.uniformLpNorm}
\difficulty{5}
Let $\Omega$ be a finite nonempty set and let $f,g:\Omega\to\mathbb{R}$. Their uniform
inner product satisfies $\langle f,g\rangle_u\leq \|f\|_{4/3,u}\,\|g\|_{4,u}$, where the
norms are the uniform $L^{4/3}$ and $L^4$ norms.
\end{theorem}

\begin{theorem}[Four third two hypercontractivity]
\lean{FABL.fourThirdTwoHypercontractivity}
\statementsource{Analysis-of-Boolean-Functions-by-Ryan-ODonnell}{
  pdf-ffbe9c7b3306-p254-b002,
  pdf-ffbe9c7b3306-p254-b003
}
\label{FABL.fourThirdTwoHypercontractivity}
\leanok
\uses{FABL.noiseOperator, FABL.twoFourHypercontractivity, FABL.twoFourNoiseRoot, FABL.uniformInner_le_uniformLpNorm_fourThird_mul_four, FABL.uniformInner_noiseOperator_left, FABL.uniformLpNorm, FABL.uniformLpNorm_nonneg, FABL.uniformLpNorm_two_sq_eq_uniformInner}
\difficulty{4}
% dataset-pipeline: stub -- prose pending build_informal_statements.py
Statement pending regeneration.
\end{theorem}

\begin{theorem}[Noise stability one third le uniform lp norm four third sq]
\lean{FABL.noiseStability_oneThird_le_uniformLpNorm_fourThird_sq}
\statementsource{Analysis-of-Boolean-Functions-by-Ryan-ODonnell}{
  pdf-ffbe9c7b3306-p254-b002,
  pdf-ffbe9c7b3306-p254-b003
}
\label{FABL.noiseStability_oneThird_le_uniformLpNorm_fourThird_sq}
\leanok
\uses{FABL.fourThirdTwoHypercontractivity, FABL.noiseOperator, FABL.noiseOperator_comp, FABL.noiseStability, FABL.noiseStability_eq_uniformInner_noiseOperator, FABL.twoFourNoiseRoot, FABL.twoFourNoiseRoot_sq, FABL.uniformInner_noiseOperator_left, FABL.uniformLpNorm, FABL.uniformLpNorm_nonneg, FABL.uniformLpNorm_two_sq_eq_uniformInner}
\difficulty{4}
% dataset-pipeline: stub -- prose pending build_informal_statements.py
Statement pending regeneration.
\end{theorem}

\begin{theorem}[Scaling of Fourier coefficients]
\lean{FABL.fourierCoeff_const_mul}
\label{FABL.fourierCoeff_const_mul}
\leanok
\uses{FABL.fourierCoeff, FABL.monomial}
\difficulty{3}
For an integer $n\ge 0$, a real-valued function $f$ on the $n$-dimensional sign cube, a
real constant $c$, and a subset $S$ of the coordinates, the Fourier coefficient of
$x\mapsto c f(x)$ at $S$ is $c$ times the Fourier coefficient of $f$ at $S$.
\end{theorem}

\begin{theorem}[Noise operator on a degree component]
\lean{FABL.noiseOperator_degreePart}
\proofsource{Analysis-of-Boolean-Functions-by-Ryan-ODonnell}{pdf-ffbe9c7b3306-p055-b003}
\label{FABL.noiseOperator_degreePart}
\leanok
\uses{FABL.degreePart, FABL.fourierCoeff_const_mul, FABL.fourierCoeff_degreePart, FABL.fourierCoeff_noiseOperator, FABL.noiseOperator, FABL.smallSet_fourier_ext}
\difficulty{4}
Let $n$ and $k$ be natural numbers, let $\rho\in\mathbb{R}$, and let
$f:\{-1,1\}^n\to\mathbb{R}$. Applying the noise operator with parameter $\rho$ to the
degree-$k$ component of $f$ gives the function $x\mapsto \rho^k$ times that degree-$k$
component evaluated at $x$.
\end{theorem}

\begin{theorem}[Scaling the fourth norm]
\lean{FABL.uniformLpNorm_four_const_mul}
\label{FABL.uniformLpNorm_four_const_mul}
\leanok
\uses{FABL.uniformLpNorm, FABL.uniformLpNorm_four_pow_eq_expect_fourth, FABL.uniformLpNorm_nonneg}
\difficulty{5}
Let $\Omega$ be a finite set, let $c\ge 0$, and let $f:\Omega\to\mathbb R$. The
uniform $L^4$ norm of $cf$ equals $c$ times the uniform $L^4$ norm of $f$.
\end{theorem}

\begin{theorem}[Fourth norm as a normalized sum]
\lean{FABL.uniformLpNorm_four_eq_sum_rpow_div_cardRoot}
\label{FABL.uniformLpNorm_four_eq_sum_rpow_div_cardRoot}
\leanok
\uses{FABL.uniformLpNorm}
\difficulty{3}
Let $\Omega$ be a finite nonempty set and let $f:\Omega\to\mathbb R$. Its uniform
$L^4$ norm equals $\bigl(\sum_{x\in\Omega}|f(x)|^4\bigr)^{1/4}/|\Omega|^{1/4}$.
\end{theorem}

\begin{theorem}[Triangle inequality for the fourth norm]
\lean{FABL.uniformLpNorm_four_add_le}
\label{FABL.uniformLpNorm_four_add_le}
\leanok
\uses{FABL.uniformLpNorm, FABL.uniformLpNorm_four_eq_sum_rpow_div_cardRoot}
\difficulty{4}
Let $\Omega$ be a finite nonempty set, and let $f,g:\Omega\to\mathbb R$. The uniform
$L^4$ norm of $f+g$ is at most the sum of the uniform $L^4$ norms of $f$ and $g$.
\end{theorem}

\begin{theorem}[Triangle inequality for finite sums]
\lean{FABL.uniformLpNorm_four_sum_le}
\label{FABL.uniformLpNorm_four_sum_le}
\leanok
\uses{FABL.uniformLpNorm, FABL.uniformLpNorm_four_add_le}
\difficulty{4}
Let $\Omega$ be a finite nonempty set, let $\iota$ be a type, let $s$ be a finite
subset of $\iota$, and let $f_i:\Omega\to\mathbb R$ for each $i\in\iota$. The
uniform $L^4$ norm of $x\mapsto\sum_{i\in s}f_i(x)$ is at most $\sum_{i\in
s}\|f_i\|_4$, where each norm is the uniform $L^4$ norm on $\Omega$.
\end{theorem}

\begin{theorem}[Two-to-four noise estimate]
\lean{FABL.homogeneous_twoFour_estimate}
\proofsource{Analysis-of-Boolean-Functions-by-Ryan-ODonnell}{
  pdf-ffbe9c7b3306-p253-b002,
  pdf-ffbe9c7b3306-p253-b003
}
\label{FABL.homogeneous_twoFour_estimate}
\leanok
\uses{FABL.degreePart, FABL.noiseOperator, FABL.noiseOperator_degreePart, FABL.twoFourHypercontractivity, FABL.twoFourNoiseRoot, FABL.twoFourNoiseRoot_nonneg, FABL.uniformLpNorm, FABL.uniformLpNorm_four_const_mul}
\difficulty{3}
Let $n$ be a nonnegative integer, let $k$ be a nonnegative integer, and let $f$ be a
real-valued function on the sign cube $\{-1,1\}^n$. Write $f_k$ for its degree-$k$
Fourier part and set $\rho=1/\sqrt{3}$. Then the uniform fourth norm of the noise
operator applied to $f_k$ equals $\rho^k$ times the uniform fourth norm of $f_k$, and
this quantity is at most the uniform second norm of $f_k$.
\end{theorem}

\begin{theorem}[Fourth-third norm of a zero-one function]
\lean{FABL.uniformLpNorm_fourThird_sq_of_zeroOne}
\label{FABL.uniformLpNorm_fourThird_sq_of_zeroOne}
\leanok
\uses{FABL.mean, FABL.uniformLpNorm}
\difficulty{4}
Let $n$ be a natural number, and let $f$ be a real-valued function on the Boolean cube
$\{-1,1\}^n$. Suppose $f$ takes only the values $0$ and $1$, and its uniform mean is
$\alpha$. Then the square of its uniform $L^{4/3}$ norm is $\alpha^{3/2}$.
\end{theorem}

\begin{theorem}[One-third noise stability bound]
\lean{FABL.smallSetExpansion_oneThird}
\proofsource{Analysis-of-Boolean-Functions-by-Ryan-ODonnell}{
  pdf-ffbe9c7b3306-p254-b002,
  pdf-ffbe9c7b3306-p254-b003,
  pdf-ffbe9c7b3306-p254-b005,
  pdf-ffbe9c7b3306-p254-b006,
  pdf-ffbe9c7b3306-p264-b001
}
\label{FABL.smallSetExpansion_oneThird}
\leanok
\uses{FABL.mean, FABL.noiseStability, FABL.noiseStability_oneThird_le_uniformLpNorm_fourThird_sq, FABL.uniformLpNorm, FABL.uniformLpNorm_fourThird_sq_of_zeroOne}
\difficulty{1}
Let $n$ be a natural number, let $f:\{-1,1\}^n\to\mathbb R$, and let $\alpha\in\mathbb
R$. Suppose $f(x)\in\{0,1\}$ for every $x$ and its uniform mean on the cube is $\alpha$.
Choose $X$ uniformly from the cube, and form $Y$ by independently keeping each
coordinate of $X$ with probability $2/3$ and flipping it with probability $1/3$. Then
$\mathbb E[f(X)f(Y)]\leq \alpha^{3/2}$.
\end{theorem}

\begin{theorem}[Small-set expansion event bound]
\lean{FABL.smallSetExpansion_oneThird_event}
\proofsource{Analysis-of-Boolean-Functions-by-Ryan-ODonnell}{
  pdf-ffbe9c7b3306-p254-b002,
  pdf-ffbe9c7b3306-p254-b003,
  pdf-ffbe9c7b3306-p254-b005,
  pdf-ffbe9c7b3306-p254-b006
}
\label{FABL.smallSetExpansion_oneThird_event}
\leanok
\uses{FABL.correlatedPairPMF, FABL.mean, FABL.noiseStability, FABL.pmfEventProbability, FABL.pmfExpectation, FABL.smallSetExpansion_oneThird}
\difficulty{4}
Let $n$ be a nonnegative integer, let $f:\{-1,1\}^n\to\mathbb R$ take only the values
$0$ and $1$, and suppose its uniform mean is $\alpha$. Draw $(X,Y)$ by choosing $X$
uniformly from the sign cube and then independently, in each coordinate, keeping its
sign with probability $2/3$ and flipping it with probability $1/3$. Then $\mathbb
P(f(X)=f(Y)=1)\le\alpha^{3/2}$.
\end{theorem}

\begin{theorem}[Three-halves power identity]
\lean{FABL.rpow_three_halves_eq_mul_sqrt}
\label{FABL.rpow_three_halves_eq_mul_sqrt}
\leanok
\difficulty{3}
For every real number $a>0$, $a^{3/2}=a\sqrt a$.
\end{theorem}

\begin{theorem}[Conditional small-set expansion bound]
\lean{FABL.smallSetExpansion_oneThird_conditional}
\proofsource{Analysis-of-Boolean-Functions-by-Ryan-ODonnell}{
  pdf-ffbe9c7b3306-p254-b005,
  pdf-ffbe9c7b3306-p264-b001
}
\label{FABL.smallSetExpansion_oneThird_conditional}
\leanok
\uses{FABL.correlatedPairPMF, FABL.mean, FABL.pmfEventProbability, FABL.rpow_three_halves_eq_mul_sqrt, FABL.smallSetExpansion_oneThird_event}
\difficulty{3}
Let $n$ be a nonnegative integer, let $f:\{-1,1\}^n\to\mathbb R$ take only the values
$0$ and $1$, and suppose its uniform mean is $\alpha>0$. Draw $(X,Y)$ by choosing $X$
uniformly from the sign cube and then independently, in each coordinate, keeping its
sign with probability $2/3$ and flipping it with probability $1/3$. Then $\mathbb
P(f(X)=f(Y)=1)/\alpha\le\sqrt{\alpha}$.
\end{theorem}

\begin{theorem}[Stable influence one third le boolean influence mul sqrt]
\lean{FABL.stableInfluence_oneThird_le_booleanInfluence_mul_sqrt}
\statementsource{Analysis-of-Boolean-Functions-by-Ryan-ODonnell}{
  pdf-ffbe9c7b3306-p255-b004,
  pdf-ffbe9c7b3306-p255-b005
}
\label{FABL.stableInfluence_oneThird_le_booleanInfluence_mul_sqrt}
\leanok
\uses{FABL.BooleanFunction, FABL.BooleanFunction.toReal, FABL.booleanInfluence, FABL.booleanInfluence_eq_influence_toReal, FABL.influence_nonneg, FABL.stableInfluence, FABL.stableInfluence_nonneg, FABL.stableInfluence_one_third_sq_le_booleanInfluence_cube}
\difficulty{4}
% dataset-pipeline: stub -- prose pending build_informal_statements.py
Statement pending regeneration.
\end{theorem}

\begin{theorem}[Stable influence one third le boolean influence rpow]
\lean{FABL.stableInfluence_oneThird_le_booleanInfluence_rpow}
\statementsource{Analysis-of-Boolean-Functions-by-Ryan-ODonnell}{
  pdf-ffbe9c7b3306-p255-b004,
  pdf-ffbe9c7b3306-p255-b005,
  pdf-ffbe9c7b3306-p264-b003
}
\label{FABL.stableInfluence_oneThird_le_booleanInfluence_rpow}
\leanok
\uses{FABL.BooleanFunction, FABL.BooleanFunction.toReal, FABL.booleanInfluence, FABL.booleanInfluence_eq_influence_toReal, FABL.influence_nonneg, FABL.rpow_three_halves_eq_mul_sqrt, FABL.stableInfluence, FABL.stableInfluence_oneThird_le_booleanInfluence_mul_sqrt}
\difficulty{4}
% dataset-pipeline: stub -- prose pending build_informal_statements.py
Statement pending regeneration.
\end{theorem}

\begin{theorem}[Low-degree Fourier weight and noise stability]
\lean{FABL.oneThird_pow_mul_fourierWeightAtMost_le_noiseStability}
\label{FABL.oneThird_pow_mul_fourierWeightAtMost_le_noiseStability}
\leanok
\uses{FABL.fourierCoeff, FABL.fourierWeight, FABL.fourierWeightAtMost, FABL.noiseStability, FABL.noiseStability_eq_sum_rho_pow_mul_sq_fourierCoeff}
\difficulty{5}
Let $n$ and $k$ be natural numbers, and let $f$ be a real-valued function on the
$n$-dimensional sign cube. Then $3^{-k}$ times the sum of the squared Fourier
coefficients of $f$ over subsets of coordinates of size at most $k$ is at most the noise
stability of $f$ at correlation $1/3$.
\end{theorem}

\begin{theorem}[Low-Degree Fourier Weight Bound]
\lean{FABL.fourierWeightAtMost_le_three_pow_mul_rpow}
\label{FABL.fourierWeightAtMost_le_three_pow_mul_rpow}
\leanok
\uses{FABL.fourierWeightAtMost, FABL.mean, FABL.oneThird_pow_mul_fourierWeightAtMost_le_noiseStability, FABL.smallSetExpansion_oneThird}
\difficulty{3}
Let $k$ be a natural number, let $f$ be a real-valued function on the $n$-dimensional
sign cube, and let $\alpha\in\bbr$. Suppose $f$ takes only the values $0$ and $1$ and
its mean is $\alpha$. Then the sum of its squared Fourier coefficients over subsets of
size at most $k$ is at most $3^k\alpha^{3/2}$.
\end{theorem}

\begin{definition}[Stable hypercube edge weight]
\lean{FABL.stableHypercubeEdgeWeight}
\statementsource{Analysis-of-Boolean-Functions-by-Ryan-ODonnell}{pdf-ffbe9c7b3306-p254-b008}
\label{FABL.stableHypercubeEdgeWeight}
\leanok
\uses{FABL.correlatedPairPMF}
Let $n$ be a nonnegative integer, let $\rho\in[-1,1]$, and let $x,y\in\{-1,1\}^n$. Their
stable edge weight is the real-valued mass assigned to the pair $(x,y)$ by the
correlated-pair probability mass function with correlation parameter $\rho$.
\end{definition}

\begin{definition}[Noisy hypercube edge weight]
\lean{FABL.noisyHypercubeEdgeWeight}
\statementsource{Analysis-of-Boolean-Functions-by-Ryan-ODonnell}{pdf-ffbe9c7b3306-p254-b008}
\label{FABL.noisyHypercubeEdgeWeight}
\leanok
\uses{FABL.one_sub_two_mul_mem_Icc, FABL.stableHypercubeEdgeWeight}
Let $n$ be a nonnegative integer, let $\delta\in[0,1]$, and let $x,y\in\{-1,1\}^n$.
Their noisy edge weight is their stable edge weight at correlation $1-2\delta$.
\end{definition}

\begin{theorem}[Nonnegativity of the edge weight]
\lean{FABL.stableHypercubeEdgeWeight_nonneg}
\proofsource{Analysis-of-Boolean-Functions-by-Ryan-ODonnell}{pdf-ffbe9c7b3306-p255-b002}
\label{FABL.stableHypercubeEdgeWeight_nonneg}
\leanok
\uses{FABL.stableHypercubeEdgeWeight}
\difficulty{1}
Let $n$ be a natural number, let $-1\leq \rho\leq 1$, and let $x,y\in\{−1,1\}^n$. The
stable hypercube edge weight from $x$ to $y$ is nonnegative.
\end{theorem}

\begin{theorem}[Point mass of the correlated pair]
\lean{FABL.stableHypercubeEdgeWeight_eq_uniform_mul_transition}
\proofsource{Analysis-of-Boolean-Functions-by-Ryan-ODonnell}{
  pdf-ffbe9c7b3306-p254-b008,
  pdf-ffbe9c7b3306-p255-b002
}
\label{FABL.stableHypercubeEdgeWeight_eq_uniform_mul_transition}
\leanok
\uses{FABL.correlatedPairPMF_apply, FABL.noiseKernel, FABL.stableHypercubeEdgeWeight, FABL.uniformPMF}
\difficulty{3}
Let $n$ be a natural number, let $-1\leq \rho\leq 1$, and let $x,y\in\{−1,1\}^n$. The
probability assigned to the pair $(x,y)$ by the correlated-pair distribution equals the
uniform probability of $x$ multiplied by the probability that the noise kernel at $x$
assigns to $y$.
\end{theorem}

\begin{theorem}[Total correlated-pair weight]
\lean{FABL.sum_stableHypercubeEdgeWeight}
\proofsource{Analysis-of-Boolean-Functions-by-Ryan-ODonnell}{pdf-ffbe9c7b3306-p255-b002}
\label{FABL.sum_stableHypercubeEdgeWeight}
\leanok
\uses{FABL.correlatedPairPMF, FABL.pmfExpectation, FABL.pmfExpectation_const_one, FABL.stableHypercubeEdgeWeight}
\difficulty{1}
Let $n$ be a natural number and let $\rho\in[-1,1]$. The sum of the stable hypercube
edge weights over all ordered pairs of points in the $n$-dimensional sign cube is $1$.
\end{theorem}

\begin{theorem}[Symmetry of the correlated pair weight]
\lean{FABL.stableHypercubeEdgeWeight_comm}
\proofsource{Analysis-of-Boolean-Functions-by-Ryan-ODonnell}{
  pdf-ffbe9c7b3306-p254-b008,
  pdf-ffbe9c7b3306-p255-b002
}
\label{FABL.stableHypercubeEdgeWeight_comm}
\leanok
\uses{FABL.correlatedPairPMF_apply_swap, FABL.stableHypercubeEdgeWeight}
\difficulty{3}
Let $\rho\in[-1,1]$, and let $x,y$ be points of the $n$-dimensional sign cube. Under the
joint distribution formed by choosing a uniform point and then independently retaining
or flipping each coordinate with retention probability $(1+\rho)/2$, the probability
assigned to the ordered pair $(x,y)$ is unchanged when $x$ and $y$ are interchanged.
\end{theorem}

\begin{definition}[Undirected stable edge weight]
\lean{FABL.stableHypercubeUndirectedEdgeWeight}
\label{FABL.stableHypercubeUndirectedEdgeWeight}
\leanok
\uses{FABL.stableHypercubeEdgeWeight}
Let $n$ be a natural number, let $\rho \in [-1,1]$, and let $x,y \in \{−1,1\}^n$. Define
the undirected stable edge weight of $x$ and $y$ at parameter $\rho$ to be the sum of
the stable edge weight from $x$ to $y$ and the stable edge weight from $y$ to $x$, using
parameter $\rho$ in both directions.
\end{definition}

\begin{theorem}[Undirected edge weight identity]
\lean{FABL.stableHypercubeUndirectedEdgeWeight_eq_two_mul}
\label{FABL.stableHypercubeUndirectedEdgeWeight_eq_two_mul}
\leanok
\uses{FABL.stableHypercubeEdgeWeight, FABL.stableHypercubeEdgeWeight_comm, FABL.stableHypercubeUndirectedEdgeWeight}
\difficulty{3}
Let $n$ be a natural number, let $-1\leq \rho\leq 1$, and let $x,y\in\{−1,1\}^n$. The
undirected stable hypercube edge weight between $x$ and $y$ equals twice the directed
stable hypercube edge weight from $x$ to $y$.
\end{theorem}

\begin{theorem}[Outgoing correlated-pair weight]
\lean{FABL.sum_stableHypercubeEdgeWeight_outgoing}
\proofsource{Analysis-of-Boolean-Functions-by-Ryan-ODonnell}{pdf-ffbe9c7b3306-p255-b002}
\label{FABL.sum_stableHypercubeEdgeWeight_outgoing}
\leanok
\uses{FABL.Sign, FABL.correlatedPairPMF_apply, FABL.noiseKernel, FABL.pmfExpectation, FABL.pmfExpectation_const_one, FABL.stableHypercubeEdgeWeight, FABL.uniformPMF}
\difficulty{5}
Let $n$ be a natural number, let $\rho\in[-1,1]$, and fix a point $x$ in the
$n$-dimensional sign cube. The sum of the stable hypercube edge weights over all points
$y$ paired with this fixed $x$ is $2^{-n}$.
\end{theorem}

\begin{theorem}[Incoming correlated-pair weight]
\lean{FABL.sum_stableHypercubeEdgeWeight_incoming}
\proofsource{Analysis-of-Boolean-Functions-by-Ryan-ODonnell}{pdf-ffbe9c7b3306-p255-b002}
\label{FABL.sum_stableHypercubeEdgeWeight_incoming}
\leanok
\uses{FABL.stableHypercubeEdgeWeight, FABL.stableHypercubeEdgeWeight_comm, FABL.sum_stableHypercubeEdgeWeight_outgoing}
\difficulty{4}
Let $n$ be a natural number, let $\rho\in[-1,1]$, and fix a point $y$ in the
$n$-dimensional sign cube. The sum of the stable hypercube edge weights over all points
$x$ paired with this fixed $y$ is $2^{-n}$.
\end{theorem}

\begin{definition}[Indicator of the all-positive point]
\lean{FABL.pointSubcubeIndicator}
\label{FABL.pointSubcubeIndicator}
\leanok
\uses{FABL.Sign}
For a natural number $k$, define a function on $\{−1,1\}^k$ that equals $1$ at the point
whose every coordinate is $1$ and equals $0$ at every other point.
\end{definition}

\begin{theorem}[Mean of a point indicator]
\lean{FABL.mean_pointSubcubeIndicator}
\label{FABL.mean_pointSubcubeIndicator}
\leanok
\uses{FABL.Sign, FABL.mean, FABL.pointSubcubeIndicator}
\difficulty{3}
For every nonnegative integer $k$, the function on the $k$-dimensional sign cube that is
$1$ at the all-$+1$ point and $0$ elsewhere has mean $2^{-k}$.
\end{theorem}

\begin{theorem}[Stay probability at a point subcube]
\lean{FABL.pointSubcube_oneThird_stayProbability}
\proofsource{Analysis-of-Boolean-Functions-by-Ryan-ODonnell}{pdf-ffbe9c7b3306-p254-b007}
\label{FABL.pointSubcube_oneThird_stayProbability}
\leanok
\uses{FABL.Sign, FABL.coordinateNoisePMF, FABL.correlationKeepProbability, FABL.independentProductPMF_apply, FABL.noiseKernel, FABL.pmfExpectation, FABL.pointSubcubeIndicator}
\difficulty{5}
For every nonnegative integer $k$, start at the all-$+1$ point of the $k$-dimensional
sign cube and apply the noise kernel with correlation parameter $1/3$. The expected
value of the indicator of that point is $(2/3)^k$.
\end{theorem}

\begin{definition}[Tensor power of a function]
\lean{FABL.pmfTensorPower}
\label{FABL.pmfTensorPower}
\leanok
For a function $g$ from a type $\Omega$ to the real numbers and a tuple
$x=(x_1,\ldots,x_d)\in\Omega^d$, define its tensor power at $x$ to be $\prod_{i=1}^d
g(x_i)$.
\end{definition}

\begin{theorem}[Moment of an independent tensor power]
\lean{FABL.exercise9_6_tensorPower_moment}
\label{FABL.exercise9_6_tensorPower_moment}
\leanok
\uses{FABL.independentProductPMF, FABL.pmfExpectation, FABL.pmfExpectation_independentProductPMF_prod, FABL.pmfTensorPower}
\difficulty{3}
Let $\Omega$ be a finite set, let $p$ be a probability mass function on $\Omega$, let
$d,m$ be natural numbers, and let $g:\Omega\to\mathbb R$. The expectation of the $m$th
power of the product $\prod_{i=1}^d g(x_i)$ under $d$ independent samples from $p$
equals $\bigl(\mathbb E_{x\sim p}[g(x)^m]\bigr)^d$.
\end{theorem}

\begin{theorem}[Exercise9 6 scaled noise decomposition]
\lean{FABL.exercise9_6_scaledNoise_decomposition}
\statementsource{Analysis-of-Boolean-Functions-by-Ryan-ODonnell}{pdf-ffbe9c7b3306-p271-b005}
\label{FABL.exercise9_6_scaledNoise_decomposition}
\leanok
\uses{FABL.degreePart, FABL.noiseOperator, FABL.noiseOperator_degreePart, FABL.noiseOperator_eq_sum_degreePart, FABL.twoFourNoiseRoot}
\difficulty{3}
% dataset-pipeline: stub -- prose pending build_informal_statements.py
Statement pending regeneration.
\end{theorem}

\begin{theorem}[Scaled noise bound by degree parts]
\lean{FABL.exercise9_6_scaledNoise_norm_le_sum}
\proofsource{Analysis-of-Boolean-Functions-by-Ryan-ODonnell}{pdf-ffbe9c7b3306-p271-b005}
\label{FABL.exercise9_6_scaledNoise_norm_le_sum}
\leanok
\uses{FABL.degreePart, FABL.exercise9_6_scaledNoise_decomposition, FABL.noiseOperator, FABL.twoFourHypercontractivity, FABL.twoFourNoiseRoot, FABL.uniformLpNorm, FABL.uniformLpNorm_four_const_mul, FABL.uniformLpNorm_four_sum_le}
\difficulty{4}
Let $n$ be a natural number, let $0<\delta\leq 1$, and let $f:\{\pm1\}^n\to\mathbb R$.
Applying the noise operator with parameter $(1-\delta)/\sqrt{3}$ satisfies
$\|T_{(1-\delta)/\sqrt{3}}f\|_4\leq\sum_{k=0}^{n}(1-\delta)^k\|f^{=k}\|_2$, where
$f^{=k}$ is the degree-$k$ Fourier part of $f$ and the norms use uniform measure on the
cube.
\end{theorem}

\begin{theorem}[Geometric square-sum bound]
\lean{FABL.exercise9_6_geometric_square_sum_le}
\label{FABL.exercise9_6_geometric_square_sum_le}
\leanok
\difficulty{5}
For every real number $\delta$ with $0<\delta\leq 1$ and every natural number $N$, the
finite geometric sum satisfies $\sum_{k=0}^{N-1}(1-\delta)^{2k}\leq 1/\delta^2$.
\end{theorem}

\begin{theorem}[Weighted degree norms bound]
\lean{FABL.exercise9_6_sum_le_inv_mul_norm}
\label{FABL.exercise9_6_sum_le_inv_mul_norm}
\leanok
\uses{FABL.degreePart, FABL.exercise9_6_geometric_square_sum_le, FABL.fourierWeightAtLevel_eq_uniformLpNorm_degreePart_sq, FABL.plancherel, FABL.sum_fourierWeightAtLevel_range, FABL.uniformLpNorm, FABL.uniformLpNorm_nonneg, FABL.uniformLpNorm_two_sq_eq_uniformInner}
\difficulty{5}
Let $n$ be a natural number, let $0<\delta\leq 1$, and let $f:\{\pm1\}^n\to\mathbb R$.
The sum of the degree-part norms weighted by $(1-\delta)^k$ satisfies
$\sum_{k=0}^{n}(1-\delta)^k\|f^{=k}\|_2\leq\delta^{-1}\|f\|_2$, with norms taken under
uniform measure on the cube.
\end{theorem}

\begin{theorem}[Scaled noise operator norm bound]
\lean{FABL.exercise9_6_scaledNoise_norm_le_inv}
\proofsource{Analysis-of-Boolean-Functions-by-Ryan-ODonnell}{pdf-ffbe9c7b3306-p271-b005}
\label{FABL.exercise9_6_scaledNoise_norm_le_inv}
\leanok
\uses{FABL.exercise9_6_scaledNoise_norm_le_sum, FABL.exercise9_6_sum_le_inv_mul_norm, FABL.noiseOperator, FABL.twoFourNoiseRoot, FABL.uniformLpNorm}
\difficulty{1}
Let $n$ be a natural number, let $0<\delta\leq 1$, and let $f$ be a real-valued function
on the $n$-dimensional Boolean cube $\{\pm1\}^n$. Applying the noise operator with
parameter $(1-\delta)/\sqrt{3}$ gives $\|T_{(1-\delta)/\sqrt{3}}f\|_4\leq
\delta^{-1}\|f\|_2$, where the norms are taken with respect to the uniform measure on
the cube.
\end{theorem}

\begin{theorem}[Positivity of the noise operator]
\lean{FABL.noiseOperator_nonneg}
\proofsource{Analysis-of-Boolean-Functions-by-Ryan-ODonnell}{pdf-ffbe9c7b3306-p063-b010}
\label{FABL.noiseOperator_nonneg}
\leanok
\uses{FABL.noiseKernel, FABL.noiseOperator, FABL.noiseOperator_apply_eq_pmfExpectation, FABL.pmfExpectation_nonneg}
\difficulty{1}
Let $n$ be a nonnegative integer, let $\rho\in[-1,1]$, and let
$f\colon\{−1,1\}^n\to\mathbb R$ satisfy $f(y)\ge0$ for every $y$. Then the noise
operator with parameter $\rho$ satisfies $(T_\rho f)(x)\ge0$ for every $x\in\{−1,1\}^n$.
\end{theorem}

\begin{theorem}[Absolute value inequality for the noise operator]
\lean{FABL.abs_noiseOperator_le_noiseOperator_abs}
\proofsource{Analysis-of-Boolean-Functions-by-Ryan-ODonnell}{pdf-ffbe9c7b3306-p063-b013}
\label{FABL.abs_noiseOperator_le_noiseOperator_abs}
\leanok
\uses{FABL.noiseKernel, FABL.noiseOperator, FABL.noiseOperator_apply_eq_pmfExpectation, FABL.pmfExpectation}
\difficulty{3}
Let $n\in\mathbb N$, let $\rho\in[-1,1]$, let $f\colon\{-1,1\}^n\to\mathbb R$, and let
$x\in\{-1,1\}^n$. Then $|T_\rho f(x)|\leq T_\rho(|f|)(x)$, where $T_\rho$ is the
operator that multiplies each Walsh coefficient indexed by a set $S$ by $\rho^{|S|}$.
\end{theorem}

\begin{theorem}[Fourth norm monotonicity]
\lean{FABL.uniformLpNorm_four_mono_abs}
\label{FABL.uniformLpNorm_four_mono_abs}
\leanok
\uses{FABL.uniformLpNorm, FABL.uniformLpNorm_four_pow_eq_expect_fourth, FABL.uniformLpNorm_nonneg}
\difficulty{4}
Let $\Omega$ be a finite set, and let $f,g:\Omega\to\mathbb R$. If $|f(x)|\le g(x)$
for every $x\in\Omega$, then the uniform $L^4$ norm of $f$ is at most that of $g$.
\end{theorem}

\begin{theorem}[Two-to-four inequality from nonnegative functions]
\lean{FABL.exercise9_7_twoFour_nonnegative_reduction}
\proofsource{Analysis-of-Boolean-Functions-by-Ryan-ODonnell}{
  pdf-ffbe9c7b3306-p253-b002,
  pdf-ffbe9c7b3306-p253-b003,
  pdf-ffbe9c7b3306-p271-b006
}
\label{FABL.exercise9_7_twoFour_nonnegative_reduction}
\leanok
\uses{FABL.abs_noiseOperator_le_noiseOperator_abs, FABL.noiseOperator, FABL.twoFourNoiseRoot, FABL.twoFourNoiseRoot_mem_Icc, FABL.uniformLpNorm, FABL.uniformLpNorm_four_mono_abs, FABL.uniformLpNorm_nonneg, FABL.uniformLpNorm_two_sq_eq_expect_sq}
\difficulty{4}
Let $n$ be a natural number, and suppose that every pointwise nonnegative function
$g:\{\pm1\}^n\to\mathbb R$ satisfies $\|T_{1/\sqrt{3}}g\|_4\leq\|g\|_2$. Then every
real-valued function $f$ on $\{\pm1\}^n$ satisfies $\|T_{1/\sqrt{3}}f\|_4\leq\|f\|_2$,
with norms taken with respect to uniform measure on the cube.
\end{theorem}

\begin{theorem}[Two-to-four noise norm inequality]
\lean{FABL.exercise9_6_twoFour_conclusion}
\proofsource{Analysis-of-Boolean-Functions-by-Ryan-ODonnell}{
  pdf-ffbe9c7b3306-p253-b002,
  pdf-ffbe9c7b3306-p253-b003
}
\label{FABL.exercise9_6_twoFour_conclusion}
\leanok
\uses{FABL.noiseOperator, FABL.twoFourHypercontractivity, FABL.twoFourNoiseRoot, FABL.uniformLpNorm}
\difficulty{1}
For every natural number $n$ and every real-valued function $f$ on $\{\pm1\}^n$, the
noise operator with parameter $1/\sqrt{3}$ satisfies $\|T_{1/\sqrt{3}}f\|_4\leq\|f\|_2$,
where both norms are taken with respect to uniform measure on the cube.
\end{theorem}
